%-------------------------------------------------------------------------------
%	SECTION TITLE
%-------------------------------------------------------------------------------
\cvsection{Eğitim}


%-------------------------------------------------------------------------------
%	CONTENT
%-------------------------------------------------------------------------------
\begin{cventries}

%---------------------------------------------------------
  \cventry
    {ELEKTRİK-ELEKTRONİK MÜHENDİSLİĞİ YÜKSEK LİSANS} % Degree
    {İzmir Katip Çelebi Üniversitesi} % Institution
    {İzmir, Türkiye} % Location
    {Eyl 2019 - Haz 2022} % Date(s)
    {
      \begin{cvitems} % Description(s) bullet points
        \textbf{Genel Not Ortalaması}: 4.00 üzerinden 3,64 (\%100 İngilizce Bölüm).\\
        \textbf{Dersler}: İstatistiksel Proses Kontrol | Veri Toplama ve Kontrol | Uygulamalı Makine Öğrenmesi | Yapay Sinir Ağları.\\
        \textbf{Tez}: Poincaré Grafiği ve İstatistiksel Analize Dayalı ML ile Tennessee Eastman Sürecinin Hata Tespiti ve Teşhisi.
      \end{cvitems}
    }

%---------------------------------------------------------
  \cventry
    {ELEKTRİK-ELEKTRONİK MÜHENDİSLİĞİ LİSANS} % Degree
    {İzmir Katip Çelebi Üniversitesi} % Institution
    {İzmir, Türkiye} % Location
    {Eyl 2014 - Tem 2019} % Date(s)
    {
      \begin{cvitems} % Description(s) bullet points
        \textbf{Genel Not Ortalaması}: 4.00 üzerinden 3,14 (\%100 İngilizce Bölüm).\\
        \textbf{Dersler}: Proses Kontrol ve Enstrümantasyon | Mühendislikte Optimizasyon | Endüstriyel Otomasyon | Sinyaller ve Sistemler.\\
        \textbf{Tez}: PLC Tabanlı İnsan-Makine Arayüzü Kontrolü ile Silindirik Kaynak ve Parlatma Makinesi Tasarımı.
      \end{cvitems}
    }

%---------------------------------------------------------
\end{cventries}
